\chapter{椭圆回路与双环网络：几何传播、耗散学习和路由控制}
\label{chp:cognitive_elliptic_resonator}

本章把椭圆看作一个可计算的低维试验台，而不是关于心智本体的缩微证明。它允许我们在同一模型中区分坐标与弧长、传播与干涉、活动衰减与信息遗失、参数学习与拓扑编辑，并观察这些过程怎样改变任务相关的状态流。原稿中的“苹果—红色”双路径、反复思虑、工作记忆、条件反射、知行切换和自我模型仍被保留，但每一种现象都必须经由明确的对象、单位、假设和可反驳观测来连接。HSF-HD在这里提供智能的一般状态动力学语言；椭圆与双环只是该语言的一组玩具模型，AgentOS则是可替换的工程实例。我们不再宣称几何必然产生心灵，而是追问表示空间、传播接口、耗散机制和结构更新各自贡献了什么，并为复杂系统建立可审计的推导基线。

\section{椭圆回路的状态空间、度量与传播}
\label{sec:elliptic_geometry_metric_curvature}

令状态载体 $\mathcal M$ 为嵌入欧氏平面的一条闭合椭圆曲线，参数 $\theta\in[0,2\pi)$，半轴 $a>b>0$，二者单位为抽象弧长单位 au，而不是米或焦耳。嵌入映射及其诱导度量为
\begin{equation}
 \mathbf r(\theta)=(a\cos\theta,b\sin\theta),\qquad
 g_{\theta\theta}=a^2\sin^2\theta+b^2\cos^2\theta .
\end{equation}
弧长元 $ds=\sqrt{g_{\theta\theta}}d\theta$，总长 $L=4aE(e)$，其中 $e=\sqrt{1-b^2/a^2}$，$E$ 是第二类完全椭圆积分。以上是定义或由定义得到的数学结果。把 $\theta$ 当作“逻辑跨度”而把 $ds$ 当作“物理能量”则没有依据：重新参数化就能改变单位角度对应的弧长，却不会改变同一条路径。可比较的量应是坐标不变的弧长、沿弧长的传播时间，或另行定义的任务代价，不能把三者混写。

原稿使用的 $\kappa(\theta)=ab(g_{\theta\theta})^{-3/2}$ 是平面嵌入曲线的外在曲率。就一维黎曼流形自身而言，局部内蕴曲率为零；仅凭曲线内部距离，不能由上式推得某处存在内在“摩擦”。本章把曲率保留为可选结构特征，例如模型可以令局部阻尼 $c(s)$ 依赖 $\kappa(s)$，但这必须写成工程假设 $c(s)=c_0+c_1f(\kappa(s))$，并通过消融实验估计 $c_1$。若 $c_1=0$ 的模型同样解释数据，曲率解释就不成立。“直觉跃迁”和“逻辑泥沼”应测为到达时间、错误率、重规划次数或计算开销；实际焦耳能耗必须由硬件功率计或可信能耗模型取得。

为描述任务，令分布状态 $z(s,t)\in\mathbb R^d$，其中 $d$ 为表示维数，$t$ 以秒计；局部结构网络另记为图 $G=(V,E)$。二者通过任务相关身份、绑定、耦合与读出接口联合，而不是把图节点等同于脑区。以“苹果”与“红色”为例，两个锚点可绑定语音和视觉证据，但是否缩短反应时间取决于传播算子、读出阈值和训练历史。若把椭圆改为同周长圆而性能不变，嵌入形状就只是可视化；若保持形状而改变阻尼后性能变化，才有证据把差异归于动力学参数。

在环上取两个锚点 $v_0$ 与 $v_\pi$，分别位于 $s=0$ 和 $s=L/2$。锚点不是把拓扑撕裂的物理钉子，而是边界或接口条件的选择。若在两点切开，环可表示为两条长度 $L/2$ 的边；若不切开，它仍是周期域，只在两点加入源与观测。令复状态 $\psi_j(s,t)\in\mathbb C$ 定义在边 $j\in\{u,d\}$，采用
\begin{equation}
 i\partial_t\psi_j=-\nu\partial_s^2\psi_j+V_j\psi_j+q_j ,
\end{equation}
其中 $\nu$ 的单位为 $\mathrm{au}^2/\mathrm{s}$，$V_j$ 为 $\mathrm{s}^{-1}$。这是研究相位传播的模型假设，不宣称认知介质服从量子力学。若接口吸收状态，可取狄利克雷条件；此时 $-\nu\partial_s^2$ 在长度 $L/2$ 上有离散本征率 $\lambda_n=\nu(2\pi n/L)^2$。离散谱只来自有限区间和边界条件，不意味着概念只能以Token存在，更不意味着连续中间态消失。若锚点是守恒连接点，条件和谱都会改变。因此椭圆提供周期状态域与非均匀坐标尺，却不提供认知体验的证明。

这一模型还需要可辨识性检查。只观察终点输出时，较长路径配合较高传播速度，可能与较短路径配合较低速度产生同样到达时间；较强阻尼配合较大输入，也可能给出同样振幅。我们因此至少在一个中间锚点采样状态，并用多组输入脉冲共同估计速度、阻尼与读出增益。参数置信区间跨越零或不同参数组同样拟合时，只能报告等价模型族，不能挑选最符合隐喻的一组。数值实现以弧长网格离散，记录 $N$、$\Delta s$、$\Delta t$ 和积分器；加密网格后结论若不稳定，所谓几何效应只是离散误差。对周期域与切开域还应使用同一数据做模型选择，防止先假定“锚点撕裂”再用该假定解释输出。

局部结构与全局表示还应在失效情形中检验。删除一个锚点身份但保留附近分布状态，系统可能仍以相似度猜出答案，却失去来源归属；清空分布状态但保留图标签，系统可能知道应查哪一路，却无法形成内容。只有两项都在、并经绑定接口联合时，模型才同时具备内容与可追踪身份。这种双重干预比脑区类比更能说明机制，也为后续双路径与双环实验提供共同基准。

\section{双路径、环路偏置与顺序依赖}
\label{sec:micro_anchoring_free_evolution}
\label{sec:gauge_field_ab_effect}

从源点到目标点存在上下两条弧，这只给出“双路径”，不自动给出“双缝干涉”。只有内部状态保留复相位、两路在读出前相干合成，输出才可写成 $y=C(\psi_u+\psi_d)$，并出现交叉项 $2\mathrm{Re}(\psi_u^*\psi_d)$。若系统传播的是概率、候选集合或独立采样，正确组合可能是 $p=\alpha p_u+(1-\alpha)p_d$，不存在相位干涉。实验可以随机打乱一路相位或身份绑定：相干模型预测输出随相位系统振荡，混合模型只预测权重平均。没有这类干预，不能把两个推理链汇合时的高置信度称为干涉证据。

“顿悟”可以改写为可检验过程。系统先沿上路检索颜色属性，沿下路检索类别和情境，在目标锚点由CCM比较两路证据。若两路均支持同一候选，读出置信度上升；若一路被污染，置信度应下降或触发重规划。高置信度可能来自条件独立证据相乘、共享表征的重复计数，也可能来自相位合成，三者不能凭主观感受区分。为防止把目标下降误当作任务成功，我们同时记录传播损失、答案正确率、校准误差和时间预算。损失变小而答案错误，说明系统只是更稳定地收敛到错误吸引子。

若系统确实采用复状态，可引入连接 $A(s,t)$ 描述相位偏置。令耦合 $q$ 无量纲，$A$ 的单位为 $\mathrm{au}^{-1}$，协变导数为 $D_s=\partial_s-iqA$。局部变换可改变单点相位，而闭合环路和乐
\begin{equation}
 \Omega(t)=q\oint_{\mathcal M}A(s,t)ds
\end{equation}
在模 $2\pi$ 意义下不变。两路合并时的相对相位包含 $\Omega$，这给出严谨的环路偏置模型。这里的 $A$ 是内部表示或控制接口参数，并非电磁矢势；“认知A-B效应”只能作为结构类比，不能据此声称存在真实磁通、贝里相位或ATP暗流。

价值规范可通过控制目标进入 $A$。例如上路强调证明完备性，下路强调按时提交，控制器依据 $h(t),E,\Theta$ 生成偏置。对普通Agent，更自然的形式可能是路径代价 $J_j=\int c_jds$ 和软路由 $p_j\propto\exp(-J_j/\tau)$；对带记忆的符号系统，则可用不对易更新算子表达顺序效应。这些机制产生相似现象，却不共享同一物理本体。原稿以非零环路偏置推导精神内耗，更稳健的做法是定义单位时间重访旧状态的次数 $r_{loop}$、未决约束残差 $e_c$ 和累计计算代价 $C_{comp}$。模型假设是：某类偏置使两个路径的局部选择相互抵消时，循环次数和计算代价上升，而任务进展率下降。固定输入并改变偏置即可检验；若只改输出相位而不增加步骤或功率，就不能称为能耗增加。

周期调度器会循环检查队列却未必存在执念，无环搜索树也可能因分支爆炸消耗大量计算。这说明循环结构既非反刍的充分条件，也非必要条件。HSF-HD描述目标、约束、输入与控制参数如何共同决定轨迹；心理解释仍需与行为和报告数据对齐。AgentOS可让TCE发现重复轨迹、SPC改变目标优先级、CCM检查是否获得新证据，但这些只是可替换实现。若 $A$ 随时间变化，导数中必须保留 $\partial_tA$；若目标、结构、输入和评价度量变化，也必须进入总导数。离散实现还需记录步长、上下文版本和随机种子，避免把配置漂移误判为拓扑效应。

还需处理价值随任务阶段改变的情形。若证明任务早期鼓励探索、临近截止时强调提交，$E(t)$ 与 $A(s,t)$ 都不是常量。比较两段轨迹必须将目标版本纳入状态，不能把策略切换误认为系统突然获得一致性。可设计交叉实验：同一输入分别配以固定目标、按计划切换目标和随机切换目标，测量循环次数、约束满足率与完成质量。只有计划切换产生可重复改善，才支持控制偏置解释；若三组无差异，环路参数并未进入有效动力学。

社会或语言任务还要防止共享证据被重复计数。上下两路若来自同一消息，合并时置信度不应按独立证据相乘。系统需要保存证据谱系，CCM按来源相关性修正权重。否则所谓相长干涉只是在两个缓存中复制同一事实。反之，两路相反也未必是相消相位，可能是一条路径使用过期数据。Me中的时间戳、版本和来源签名因此属于模型的一部分，而不是实现细节。

\section{耗散、保持时间与慢结构学习}
\label{sec:dissipation_complex_spectrum}
\label{sec:dissipative_heat_and_plastic_collapse}

工作记忆先从泄漏系统描述。令 $z(s,t)\in\mathbb R^d$ 为活动状态，$D$ 的单位为 $\mathrm{au}^2/\mathrm{s}$，泄漏率 $\gamma(s,t)\ge0$ 的单位为 $\mathrm{s}^{-1}$，输入矩阵 $B$ 把事件 $u(t)$ 映射到状态空间：
\begin{equation}
 \partial_tz=D\partial_s^2z-\gamma z+Bu(t),\qquad y=Cz .
\end{equation}
在均匀泄漏、无输入且忽略扩散时，$z(t)=z(0)e^{-\gamma t}$；活动范数半衰期为 $\ln2/\gamma$，平方范数半衰期为 $\ln2/(2\gamma)$，二者不能混称“记忆寿命”。真正的任务保持时间还依赖读出阈值 $\varepsilon$，即 $T_{keep}=\gamma^{-1}\ln(\|z(0)\|/\varepsilon)$。若状态以有效速度 $v$ 传播，阈值为初值比例 $\rho$，可定义保持距离 $\xi_\rho=(v/\gamma)\ln(1/\rho)$，但只在速度近似稳定、衰减近似指数且途中不刷新时成立。

采用复状态也可写 $i\partial_t\psi=(H-i\gamma I)\psi$，但复本征值只是线性开放系统的表示，不证明遗忘服从单一指数。多时间尺度会产生指数和，检索干扰可产生非单调曲线，外部记忆Me可能在内部活动消失后恢复内容。实验必须比较单指数、多指数、幂律和干扰模型。“想到一半忘了前提”可能是路径过长，也可能是新输入覆盖、身份绑定错误或检索失败，需要冻结输入、关闭竞争记忆和改变路径长度分别识别。

对二次活动量 $Q=\tfrac12\|z\|^2$，周期或无通量边界下有
\begin{equation}
 \dot Q=-D\|\partial_sz\|^2-\int\gamma\|z\|^2ds+\langle z,Bu\rangle .
\end{equation}
它描述模型活动量收支，不是焦耳能量守恒。实际电功率需另测芯片、内存与通信。Offloading把状态写入Me时，内部 $Q$ 可下降而任务性能上升，恰好说明活动强度、信息保存和硬件能耗是不同量。高激活也可能只是噪声大，较小交叉熵亦不保证答案满足现实约束。

快速传播用 $t$ 描述，结构学习用慢索引 $\tau$。令 $w(\tau)$ 包括边权、锚点绑定、读出矩阵和正定度量密度 $m(s,\tau)$。训练目标、结构、输入分布和评价度量分别记为 $E_\tau,G_\tau,P_\tau,M_\tau$：
\begin{equation}
 \mathcal L(w,\tau)=\mathbb E_{x\sim P_\tau}\ell(F(x;w,G_\tau),E_\tau;M_\tau)+\lambda R(w),
\end{equation}
并用 $w_{k+1}=\Pi_{\mathcal W}[w_k-\eta_k\widehat\nabla_w\mathcal L_k]$ 更新。若目标、结构、输入或评价标准改变，总导数必须包含相应项，不能由某次损失下降推出任务已成功。

“看到火—感到痛”的重复配对可作为案例。初始时两锚点之间需要多步证据传播；联合出现产生预测误差，学习器提高直接边权或建立缓存，反应时下降。快状态消散与慢权重改变由不同方程描述，可由资格迹和误差信号联结，但没有误差或可塑更新时，活动会消失却不会留下结构。这直接否定“耗散废热必然凝固为记忆”。度量学习可降低频繁成功路径的算法距离，但只改变椭圆在平面中的形状、保持弧长与算子不变时，性能不会变化。快捷边还可能在情境反转后误触发，所以Boundary应保存适用域，CCM在高风险输入上重启慢路径，TCE检测漂移，Me保存可回滚版本。塑性由此是可验证的结构学习，而非能量名称的机械替换。

时间尺度分离不是自动成立的。若快速状态尚未收敛，慢更新便依据瞬时残差修改结构，系统可能追逐噪声并形成错误捷径。实现时可规定：只有CCM残差在连续 $K$ 个窗口内低于阈值，或证据累计达到预注册标准，TCE才提交结构更新；否则更新留在影子副本。若采用联合在线学习，则必须报告快慢步长比、延迟和稳定区间。比较前后模型时使用冻结测试集，并固定目标版本、输入分布和评价度量。

遗忘与外部记忆之间也存在资源交换。写入Me会增加序列化、索引和检索成本，却降低内部维持负担；是否值得取决于未来复用概率、检索延迟和错误恢复成本。总任务代价可由内部更新、外部读写和超时惩罚加权，但这些权重仍是工程尺度，不是焦耳。硬件能耗需另测。系统可能计算代价下降而电能上升，也可能信息熵下降而任务失败，任何单项指标都不能替代闭环验收。

条件反射案例还应加入去学习与再学习阶段。若火源已隔离，旧快捷边应在反证积累后降权；若风险重新出现，系统应从Me恢复证据而非从零训练。记录撤销延迟、错误暴露和恢复样本数，才能判断结构记忆是否既稳定又可修订。不可撤销的快速反应也许适合恒定环境，却不应概括为更高智能。

\section{从单环到双环：拓扑编辑、路由与自我模型}
\label{sec:topological_phase_transition_lemniscate}
\label{sec:quantum_dynamics_and_self_nucleation}

椭圆无论怎样连续拉伸或压扁，只要不剪切、不粘合且保持一一连续映射，拓扑仍是 $S^1$，第一贝蒂数 $\beta_1=1$。要得到双环楔和 $S^1\vee S^1$，必须明确执行商空间粘合、增加边或合并节点。把平面椭圆画成自交“8”字只改变嵌入图像；若交叉处两支不连通，抽象拓扑仍是单环。只有交叉点被声明为共享节点并允许换路，图的循环秩才变为 $|E|-|V|+1=2$。因此相变不是曲率积累的必然结果，而是带前置条件的离散结构编辑。

令候选操作 $T$ 把两个锚点身份合并为节点 $x$，同时保留两条独立闭路。操作前应验证绑定相容性、权限、资源预算和回滚快照；操作后计算连通分量、循环秩和接口测试。事务分为准备、影子运行、切换、验收：先冻结版本，再比较新旧图输出，随后原子更新路由表，失败则恢复旧图。度量接近不等于身份相同，更不授权改变结构。双环可把一环配置为TDCI内部推演，另一环配置为TECI环境交互，也可按语言与视觉或规划与校验划分。功能来自节点标签、边算子和权限，不来自孔洞数量；相同 $\beta_1$ 的图可以实现完全不同的Agent，单环加高维状态也可能完成内外协同。

在共享节点 $x$，若采用概率路由，令 $p(e_{out}\mid e_{in},h,E,\Theta)$ 为行随机矩阵；若采用复振幅模型，则散射矩阵 $S$ 需满足 $S^*S=I$。二者不可混用：概率守恒不产生相位干涉，幺正振幅需要相干假设。节点状态汇总任务身份、最近证据和资源预算，路由由控制策略生成。所谓选择因此是受状态和目标调制的决策，不是节点度数自动涌现的自由意志。

TDCI环可运行反事实模拟，TECI环可提交动作并接收反馈。一个周期是可追踪序列：内部环提出候选，CCM检查约束，交叉点决定继续推演或送往外部环，结果再写回Me。若两环目标一致，切换成本和冲突残差可能下降；若不一致，TCE应请求SPC重新排序。“心流”最多是低切换成本、高进展率和低错误率的经验标签，不能由两项虚构磁通互为相反数直接推出。

共享节点适合放置自我模型 $\Psi_{self}$，因为它同时读取内部承诺和外部后果。我们将其定义为对系统边界、能力、承诺和控制权的可更新表示，输出身份集合、可用动作、承诺账本和不确定度。Boundary据此检查动作是否越权以及失败由谁承担，观测再更新模型。这是工程设计；其他智能可以用分布式协议完成同一功能。局部交叉节点与全局分布表示并不等同：身份键把分布证据绑定到节点端口，耦合算子控制读写，接口生成可审计决定。删除身份键会误合并不同主体的承诺，只有键而无分布内容又无法利用上下文。由“枢纽位置”推出“自我必然成核”跨越了机制层级。

验证应注入三类故障：冻结自我模型后改变权限，交换两环身份标签，以及保持循环秩不变而改变边延迟。若路由随标签而非表面位置改变，说明绑定决定功能；若无中心节点的分布式实现也完成任务，“自我必须位于奇点”即被否定。双环收益还需与单环加捷径、外部缓存和并行检索在相同预算下比较。HSF-HD在这里描述结构状态与快速流的耦合，AgentOS提供SPC、TCE、CCM、Boundary和Me的一个实现，但它们都不是所有智能系统的唯一机制。

\section{验证协议、反例与理论边界}
\label{sec:conclusion_1d_resonator}

本章最终给出可复现实验，而不是从自动机到生命的断言。实验一比较同周长圆与椭圆：固定传播算子、锚点和数据，只改变嵌入形状；若结果不变，外在曲率没有因果作用。随后固定形状、改变度量密度或阻尼，检验到达时间与错误率。实验二比较相干合成和概率混合：对一路注入受控相位扰动，只有前者应出现周期输出。实验三估计泄漏率和保持阈值，用多种衰减模型交叉验证，并将内部保持、Me检索和硬件功率分开报告。

实验四针对慢学习。对“火—痛”配对设置训练、情境反转和恢复阶段，记录快捷路由时延、误触率、慢路径重启率和参数漂移。若训练只缩短时延却无法在反转后恢复，形成的是脆弱捷径而非一般能力。实验五比较单环加捷径与双环，在相同参数量、计算预算和数据下评估；只有双环在多项指标上稳定占优，才支持拓扑编辑的工程价值。实验六测试自我模型：改变权限、承诺和控制主体，检查系统是否正确更新责任归属，而不是以语言自述替代行为证据。

所有实验都要区分五类陈述。对象和单位属于定义；边界条件、泄漏形式和路由族属于模型假设；给定假设下的谱、半衰期与循环秩属于数学结果；顿悟、反刍、心流和自我体验属于需要数据的经验主张；把SPC、TCE、CCM、Boundary、Offloading和Me组合起来属于AgentOS工程设计。某一层成立不能自动提升为另一层成立。尤其不能由预测有效推出整体同构，不能由损失下降推出任务成功，也不能把熵、活动量、计算代价和焦耳能耗互换。

我们保留三个反例作为护栏。第一，双路径可以只有混合而无干涉；第二，耗散可以没有学习，学习也可通过外部写入而非耗散发生；第三，双环可以没有自我模型，单环或无环系统也可能实现边界与责任跟踪。因此椭圆、和乐和双环都不是智能的必要或充分条件。它们的价值在于构成最小可分析模型，使我们逐项观察几何、状态、控制和结构的作用。

数值验证也不可省略。空间网格步长 $\Delta s=L/N$、时间步长 $\Delta t$ 和慢更新步长 $\eta_k$ 必须记录；显式扩散更新需满足所选格式的稳定条件，并做网格收敛检查。在线目标变化时，前后模型须在冻结测试集上重算。结构编辑须保存版本、事务日志与回滚结果。行为指标至少包括正确率、校准误差、时延、计算预算、危险误触和恢复时间，任何单一指标改善都不足以宣告任务成功。

从全书层级看，系统论、协同论、耗散理论、复杂系统理论和认知科学提供问题来源；HSF-HD把问题组织为状态、目标、边界、耦合与多时间尺度演化；AgentOS再把其中一部分落实为接口和验证流程。这个层级避免把HSF-HD等同于某个操作系统，也避免借物理名词替代证明。我们由此得到较克制但更强的结论：智能系统能够通过可测的传播、泄漏、刷新、结构更新和路由控制，把反复经验转化为更快而仍可回滚的行为；是否产生生命或意识，超出本章模型的可证明范围，必须交由更广泛的经验研究回答。

排版和接口验证也是论证的一部分。公式中的对象、单位与假设必须在首次出现处可见，长公式不得越界，旧标签应落在语义对应位置，使历史引用仍指向相同问题。若章节在源码计数上合格却因分页只留下几行正文，仍不能通过版面审阅。审阅PDF应逐页检查半页门槛、公式断行、交叉引用和中文字体，并把页码与结论写回验证记录。

验收还必须允许失败。若圆与椭圆实验无差异，我们删除曲率因果主张；若概率混合比相干模型更好，就停止使用干涉解释；若双环不优于单环捷径，就保留较简单结构；若自我模型对越权故障无改善，则回到边界接口重新设计。理论能够规定怎样被否定，才比不可证伪的宏大比喻更有解释力。

%---chp---
